\documentclass[10pt,a4paper]{amsart}
\usepackage[margin=19mm]{geometry}
\usepackage{amsmath,amssymb,mathtools}
\usepackage[scr=rsfs,cal=euler]{mathalfa}
\usepackage{xfrac}
\usepackage[unicode,hidelinks,bookmarks=false]{hyperref}
\newcommand{\ie}{i.e.\ }
\newcommand{\cf}{cf.\ }
\newcommand{\resp}{respectively\ }
\newcommand{\Subsection}{\subsection}
\providecommand{\nobreakdash}{\nobreak\hbox{-}}
\setlength{\emergencystretch}{2em}
\multlinegap0pt
\usepackage{amssymb}
\usepackage[all]{xy}
\usepackage{mathtools}
\usepackage{braket}
\newtheorem{prop}{Proposition}
\title{Variations of the Hardy Z-Function and the Montgomery Pair Correlation Conjecture}
\author{Yochay Jerby}
\date{Source: arXiv:2511.18275v2 (CC BY 4.0)}
\begin{document}
\maketitle

\noindent\emph{Source and licence.} Variations of the Hardy Z-Function and the Montgomery Pair Correlation Conjecture. Version arXiv:2511.18275v2 (CC BY 4.0), distributed by arXiv under the Creative Commons Attribution 4.0 International licence, http://creativecommons.org/licenses/by/4.0/, as recorded in the arXiv licence metadata for this exact version and shown on the abstract page https://arxiv.org/abs/2511.18275v2. Full licence notice: \texttt{licenses/arxiv-2511.18275-CC-BY-4.0.txt}.\par\smallskip

\noindent\textit{Editorial scope.} Counted unit: the article's prop prop:breg-regular-cut (source0.tex lines 2087--2117). The article's complete original proof of it is delivered with it (source0.tex lines 2120--2350, one \texttt{proof} environment of 231 lines). No mathematics is written, repaired or re-derived: the statement and the proof are the article's own lines, in the article's own notation, and no retained line is substituted or re-flowed.  No further source-local context is needed: every cross-reference of every retained line resolves inside the two delivered spans.  Nothing was omitted from any retained span, so the package carries no omission record.  No retained line contains a citation, so the wrapper carries no bibliography and no bibliography entry.  No retained line contains a figure environment, an image-inclusion command or drawing code, so no image is delivered and the package declares no asset dependency.  Editorial formatting only, in addition: an AMS \texttt{amsart} preamble that restates the article's own declarations and macros used by the retained lines, namely the environments \texttt{prop} on separate counters, together with the standard packages those lines need.  The retained environments print in this document as Proposition 1 in order of appearance, whereas the article prints the same items with the numbering of its own full document.  The complete native source of the article as distributed in the arXiv e-print for this exact version is bundled unchanged as \texttt{sources/189641/source0.tex}.
\begin{prop}[Admissibility of $b^{\mathrm{reg}}$ for universality]
\label{prop:breg-regular-cut}
Fix $S>0$ and $\delta\in(0,1]$. Then there exist constants
$C_{S},L_{S,\delta},N_0<\infty$ such that for all
$N\ge N_0$:

\begin{enumerate}
\item[(i)] The following bound holds 
\begin{equation}\label{eq:breg-sup-cut}
\sup_{t\in[0,S]}\ \sup_{n\in I_N}\ \big|\,b^{\mathrm{reg}}_{n}(t;\delta)\,\big|
\;\le\; C_{S}\,\frac{1}{\delta h_N}\,\big(1+\log |I_N|\big).
\end{equation}

\item[(ii)] Let
\begin{equation}\label{eq:D-delta}
\mathcal D_{\delta}:=\Big\{(\widetilde X_k)_{k\in I_N}\ :\
|\widetilde X_{k+1}-\widetilde X_k|\ge \delta\,h_N\ \text{for all }k\in I_N\Big\}.
\end{equation}
For every $t\in[0,S]$ the map
$(\widetilde X_k)_{k\in I_N}\mapsto b^{\mathrm{reg}}(t;\delta)$ is $C^\infty$
on $\mathcal D_\delta$, and for all configurations
$\widetilde X,\widetilde Y\in\mathcal D_\delta$,  the following $\ell_1$ Lipschitz condition holds
\begin{equation}\label{eq:breg-Lip}
\sum_{n\in I_N}\big|\,b^{\mathrm{reg}}_{n}(\widetilde X;t;\delta)
-b^{\mathrm{reg}}_{n}(\widetilde Y;t;\delta)\,\big|
\ \le\ L_{S,\delta}\,
\frac{1+\log |I_N|}{(\delta h_N)^2}\,
\sum_{k\in I_N}|\widetilde X_k-\widetilde Y_k| .
\end{equation}
\end{enumerate}
\end{prop}

\begin{proof} 

(i) Write $b^{\mathrm{reg}}_{n}=T_{1,n}+T_{2,n}$ with
\begin{equation}
T_{1,n}(t):=\sum_{m\neq n}
\Big(1-\chi_N\!\big(\widetilde X_n-\widetilde X_m;\delta\big)\Big)
\frac{\widetilde c_{nm}(A_t)-1}{\widetilde X_n-\widetilde X_m} 
\end{equation}
\begin{equation}
T_{2,n}(t):=\sum_{m\neq n}
\Big(1-\chi_N\!\big(\widetilde X_n-\widetilde X_m;\delta\big)\Big)\, r^{\mathrm{gap}}_{nm}(t),
\end{equation}
By construction, $1-\chi_N(\widetilde X_n-\widetilde X_m;\delta)$ vanishes unless
$|\widetilde X_n-\widetilde X_m|\ge \delta h_N$, up to a harmless smooth
transition. Hence, in the relevant region 
\begin{equation}\label{eq:LB}
|\widetilde X_n-\widetilde X_{n\pm j}|
\;\ge\; \sum_{k=0}^{j-1}|\widetilde X_{n+k+1}-\widetilde X_{n+k}|
\;\ge\; j\,\delta h_N .
\end{equation}
On the fixed window $[0,S]$ the coefficients $\widetilde c_{nm}(A_t)$ are
uniformly bounded, hence
\begin{equation}
\label{eq:bound-1}
|T_{1,n}(t)|
\;\le\; C_1 \sum_{m\neq n}\frac{1-\chi_N\!\big(\widetilde X_n-\widetilde X_m;\delta\big)}
{|\widetilde X_n-\widetilde X_m|}
\;\le\; \frac{C_1}{\delta h_N}\sum_{j=1}^{|I_N|}\frac{2}{j}
\;\le\;2  C_1 \,\frac{1}{\delta h_N}\,\big(1+\log |I_N|\big).
\end{equation} 
Recall that
\begin{equation}
r^{\mathrm{gap}}_{nm}
=\widetilde c_{nm}\left[\frac{1}{\phi_{nm}(A_t)\,(X_n-X_m)}
-\frac{1}{\widetilde X_n-\widetilde X_m}\right].
\end{equation}
When $1-\chi_N\neq 0$, both denominators are $\ge c\,\delta h_N$ and
$\phi_{nm},\widetilde c_{nm}$ are uniformly bounded on $[0,S]$. Hence 
\begin{equation}
|r^{\mathrm{gap}}_{nm}(t)|\le C_{2}/|\widetilde X_n-\widetilde X_m|.
\end{equation}
Therefore,
\begin{equation}
\label{eq:bound-2}
|T_{2,n}(t)|
\;\le\; \frac{C_2}{\delta h_N}\sum_{j\ge1}\frac{2}{j}
\;\le\; C_{2}\,\frac{1}{\delta h_N}\,\big(1+\log |I_N|\big).
\end{equation}
Combining the two bounds \eqref{eq:bound-1} and \eqref{eq:bound-2} gives \eqref{eq:breg-sup-cut}.

(ii) Set
\begin{equation}
\Delta_{nm}\;:=\;\widetilde X_n-\widetilde X_m,\qquad
\chi_{nm}\;:=\;\chi_N(\Delta_{nm};\delta)
=\chi\!\Big(\frac{|\Delta_{nm}|}{\delta h_N}\Big).
\end{equation}
Note that
\begin{equation}
b^{\mathrm{reg}}_{n}(\widetilde X;t;\delta)
=\sum_{m\neq n}\Big(1-\chi_{nm}\Big)\Bigg[
      \frac{\widetilde c_{nm}(A_t)-1}{\Delta_{nm}}
      + r^{\mathrm{gap}}_{nm}(t;\widetilde X)\Bigg].
\end{equation}
By definition of $\mathcal D_\delta$,
\begin{equation} 
|\widetilde X_{k+1}-\widetilde X_k|\ge \delta h_N
\end{equation} for all $k$. Consequently, for any $n\neq m$ we have the telescoping lower bound
\begin{equation}\label{eq:LB-j}
|\Delta_{nm}|
=\Big|\sum_{j=0}^{|m-n|-1}
(\widetilde X_{n+j+1}-\widetilde X_{n+j})\Big|
\;\ge\; |m-n|\,\delta h_N .
\end{equation}
In particular, on $\mathcal D_\delta$ every denominator
$\Delta_{nm}$ in the above expression is bounded away from zero once the cutoff factor $1-\chi_{nm}$ is present.
The maps $\Delta_{nm} \mapsto 1/\Delta_{nm}$ and the cutoff $\chi$ are $C^\infty$ on $\{\Delta\neq0\}$.
Since $r^{\mathrm{gap}}_{nm}$ is a finite linear combination of such
fractions with $A_t$–dependent coefficients,
each summand is a $C^\infty$ function of $\{\widetilde X_k\}$ on
$\mathcal D_\delta$ and the sum converges absolutely. Therefore $(\widetilde X_k)\mapsto b^{\mathrm{reg}}(t;\delta)$
is $C^\infty$ on $\mathcal D_\delta$.


For the $\ell_1$ Lipschitz bound, let $\widetilde X,\widetilde Y\in\mathcal D_\delta$ and set
\begin{equation}
\widetilde Z(\theta):=\widetilde Y+\theta(\widetilde X-\widetilde Y) \in \mathcal{D}_{\delta},
\end{equation} for
$\theta\in[0,1]$, be their linear interpolation, which is well defined due to the convexity of $\mathcal{D}_{\delta}$. By the fundamental theorem of calculus and the chain rule,
\begin{equation}
b^{\mathrm{reg}}_{n}(\widetilde X;t;\delta)-b^{\mathrm{reg}}_{n}(\widetilde Y;t;\delta)
=\int_0^1\sum_{k\in I_N}
\partial_{\widetilde X_k} b^{\mathrm{reg}}_{n}(\widetilde Z(\theta);t;\delta)\,
\big(\widetilde X_k-\widetilde Y_k\big)\,d\theta .
\end{equation}
Summing over $n$ gives
\begin{equation}\label{eq:FToC}
\sum_{n\in I_N}\big|b^{\mathrm{reg}}_{n}(\widetilde X;t;\delta)
-b^{\mathrm{reg}}_{n}(\widetilde Y;t;\delta)\big|
\;\le\; \Big(\sup_{\theta\in[0,1]}
\max_{k\in I_N}\,\sum_{n\in I_N}
\big|\partial_{\widetilde X_k} b^{\mathrm{reg}}_{n}(\widetilde Z(\theta);t;\delta)\big|\Big)
\sum_{k\in I_N}|\widetilde X_k-\widetilde Y_k|.
\end{equation}
Thus it suffices to bound uniformly on $\mathcal D_\delta$ the quantity
\begin{equation}
\Lambda_{k}\;:=\;\sum_{n\in I_N}
\big|\partial_{\widetilde X_k} b^{\mathrm{reg}}_{n}(\widetilde Z;t;\delta)\big|.
\end{equation}
\noindent
Differentiating a single summand in the definition of $b^{\mathrm{reg}}_{n}$ with
respect to $\widetilde X_k$ gives, for each $m\neq n$,
\begin{multline}\label{eq:def-alpha}
\alpha_{nm}^k
:=\partial_{\widetilde X_k}\!\Big[
(1-\chi_{nm})\Big(\frac{\widetilde c_{nm}-1}{\Delta_{nm}}
+ r^{\mathrm{gap}}_{nm}\Big)\Big]=\\
=(1-\chi_{nm})\,
\partial_{\widetilde X_k}\!\Big(\frac{\widetilde c_{nm}-1}{\Delta_{nm}}
+ r^{\mathrm{gap}}_{nm}\Big)
\;-\;\chi'_{nm}\,
\partial_{\widetilde X_k}\!\Big(\frac{|\Delta_{nm}|}{\delta h_N}\Big)\,
\Big(\frac{\widetilde c_{nm}-1}{\Delta_{nm}}+ r^{\mathrm{gap}}_{nm}\Big),
\end{multline}
where $\chi'_{nm}:=\chi'\!\big(|\Delta_{nm}|/(\delta h_N)\big)$.  Note that
\begin{equation}\label{eq:derivs-Delta}
\partial_{\widetilde X_k}\Delta_{nm}=\mathbf{1}_{\{k=n\}}-\mathbf{1}_{\{k=m\}},
\qquad
\partial_{\widetilde X_k}\!\Big(\frac{|\Delta_{nm}|}{\delta h_N}\Big)
=\frac{\operatorname{sgn}(\Delta_{nm})}{\delta h_N}
\big(\mathbf{1}_{\{k=n\}}-\mathbf{1}_{\{k=m\}}\big).
\end{equation}
Set
\begin{equation}
F_{nm}(\widetilde X)
:= \frac{\widetilde c_{nm}(A_t)-1}{\Delta_{nm}}
   + r^{\mathrm{gap}}_{nm}(t).
\end{equation}
The quantities $\widetilde c_{nm}(A_t)$ and $\phi_{nm}(A_t)$ appearing in
$r^{gap}_{nm}$ depend only on the time–parameter $A_t$.
Hence $F_{nm}$ depends on the configuration $(\widetilde X_k)_{k\in I_N}$
only through the two coordinates $\widetilde X_n$ and $\widetilde X_m$, via
$\Delta_{nm}$ and
\begin{equation}
X_n-X_m=\sigma_n(A_t)\,\widetilde X_n-\sigma_m(A_t)\,\widetilde X_m .
\end{equation}
where $\big\|\nabla t_n(A_t) \big\|$, which appears in $r^{gap}_{nm}$. By the chain rule,
\begin{equation}
\partial_{\widetilde X_k}F_{nm}(\widetilde X)
=\partial_{\Delta}F_{nm}(\widetilde X)\;\partial_{\widetilde X_k}\Delta_{nm}
+\partial_{(X_n-X_m)}F_{nm}(\widetilde X)\;\partial_{\widetilde X_k}(X_n-X_m).
\end{equation}
Direct differentiation gives
\begin{equation}
\partial_{\widetilde X_k}\Delta_{nm}
=\mathbf 1_{\{k=n\}}-\mathbf 1_{\{k=m\}},
\qquad
\partial_{\widetilde X_k}(X_n-X_m)
=\sigma_n(A_t)\mathbf 1_{\{k=n\}}-\sigma_m(A_t)\mathbf 1_{\{k=m\}} .
\end{equation}
Therefore both derivatives vanish whenever $k\notin\{n,m\}$, and thus
\begin{equation}
(1-\chi_{nm})\,\partial_{\widetilde X_k}
\Big(\frac{\widetilde c_{nm}(A_t)-1}{\Delta_{nm}}+r^{\mathrm{gap}}_{nm}(t)\Big)=0
\qquad\text{if }k\notin\{n,m\}.
\end{equation}
For the cutoff term, note also that
\begin{equation}
\partial_{\widetilde X_k}\!\Big(\frac{|\Delta_{nm}|}{\delta h_N}\Big)
=\frac{\operatorname{sgn}(\Delta_{nm})}{\delta h_N}
\big(\mathbf 1_{\{k=n\}}-\mathbf 1_{\{k=m\}}\big),
\end{equation}
so the derivative of $\chi\!\left(\frac{|\Delta_{nm}|}{\delta h_N}\right)$
is zero as well when $k\notin\{n,m\}$. Consequently, $\alpha^k_{nm}=0$
unless $k\in\{n,m\}$.

\medskip\noindent
\emph{Case $k=n$.} Then every $m\neq n$ contributes. On $\mathcal D_\delta$ the
telescoping lower bound
\begin{equation}\label{eq:LB-j-again}
|\Delta_{n,n\pm j}|=\Big|\sum_{\ell=0}^{j-1}
(\widetilde X_{n+\ell+1}-\widetilde X_{n+\ell})\Big|
\ \ge\ j\,\delta h_N,\qquad j\ge 1,
\end{equation}
holds. Using that $\widetilde c_{nm}$, $r^{\mathrm{gap}}_{nm}$ and their first
derivatives are uniformly bounded on $[0,S]$, we get
\begin{equation}\label{eq:sum-alpha-n}
\sum_{m\neq n}\!|\alpha_{n m}^{\,n}|
\;\lesssim\; \sum_{j\ge1}\!\left(\frac{2}{(j\,\delta h_N)^2}
\;+\;\frac{2}{\delta h_N}\cdot
\frac{\mathbf 1_{\{j\,\delta h_N\in[\delta h_N,\,2\delta h_N]\}}}{j\,\delta h_N}\right)
\;\le\; \frac{C}{(\delta h_N)^2}\Big(1+\log |I_N|\Big).
\end{equation}

\medskip\noindent
\emph{Case $k\neq n$.} Then only the pair $(n,m)=(n,k)$ contributes. By
\eqref{eq:LB-j-again} with $j=|k-n|$,
\begin{equation}\label{eq:alpha-nk}
|\alpha_{n k}^{\,k}|
\;\lesssim\; \frac{1}{(|k-n|\,\delta h_N)^2}
\;+\;\frac{1}{\delta h_N}\cdot
\frac{\mathbf 1_{\{|k-n|\,\delta h_N\in[\delta h_N,\,2\delta h_N]\}}}{|k-n|\,\delta h_N}.
\end{equation}
Summing \eqref{eq:alpha-nk} over $n\neq k$ yields
\begin{equation}\label{eq:sum-alpha-k}
\sum_{n\neq k}\!|\alpha_{n k}^{\,k}|
\;\le\; \frac{C}{(\delta h_N)^2}\Big(1+\log |I_N|\Big).
\end{equation}

\medskip
Combining \eqref{eq:sum-alpha-n} and \eqref{eq:sum-alpha-k} and recalling
$\Lambda_k=\sum_{n\in I_N}|\partial_{\widetilde X_k}
b^{\mathrm{reg}}_{n}(\widetilde Z;t;\delta)|$ gives the uniform bound
\begin{equation}\label{eq:Lambda-k-final}
\Lambda_k
=\sum_{n\in I_N}\big|\partial_{\widetilde X_k} b^{\mathrm{reg}}_{n}(\widetilde Z;t;\delta)\big|
\;\le\; \frac{C_{S,\delta}}{(\delta h_N)^2}\Big(1+\log |I_N|\Big),
\qquad \text{for all }\ \widetilde Z\in\mathcal D_\delta.
\end{equation}
Substituting \eqref{eq:Lambda-k-final} into \eqref{eq:FToC} completes the proof
of the $\ell_1$–Lipschitz estimate.
 Inserting \eqref{eq:Lambda-k-final} into \eqref{eq:FToC} we conclude the
\(\ell_1\)–Lipschitz estimate
\begin{equation}
\sum_{n\in I_N}\big|b^{\mathrm{reg}}_{n}(\widetilde X;t;\delta)
-b^{\mathrm{reg}}_{n}(\widetilde Y;t;\delta)\big|
\ \le\ L_{S,\delta}\,
\frac{1+\log |I_N|}{(\delta h_N)^2}\,
\sum_{k\in I_N}|\widetilde X_k-\widetilde Y_k|,
\end{equation}
as claimed.
\end{proof}

\end{document}
